\chapter{Estudo de Caso: Patch-Hub}

\section{Visão Geral do Projeto}

Este trabalho examina a reescrita do Patch-Hub, uma ferramenta de código aberto desenvolvida sob a organização KWorkflow (kw), principalmente por estudantes da USP. O Patch-Hub é uma interface baseada em terminal para navegar e interagir com patches em lore.kernel.org, suportando o kernel Linux e projetos relacionados. Ele agiliza o fluxo de trabalho tradicional centrado em listas de discussão usado no desenvolvimento do kernel.

O Patch-Hub pertence ao projeto mais amplo KWorkflow, um Sistema de Fluxo de Trabalho de Automação de Desenvolvedor que reduz a sobrecarga de configuração do kernel e fornece ferramentas para tarefas do dia a dia. Reescrever o Patch-Hub usando o Modelo de Atores visa tanto melhorar a ferramenta quanto demonstrar os benefícios de arquiteturas baseadas em atores para aplicações centralizadas em Rust.

Este projeto é um caso de estudo adequado porque:

\begin{itemize}
    \item Tem um escopo e conjunto de funcionalidades claros
    \item Apresenta desafios arquiteturais bem adequados ao Modelo de Atores
    \item É mantido ativamente e usado pela comunidade do kernel
    \item Permite medição concreta do impacto arquitetural
\end{itemize}

\section{Arquitetura Atual}

A implementação atual do Patch-Hub segue uma arquitetura principalmente single-threaded e monolítica com os seguintes módulos principais:

\begin{itemize}
    \item \textbf{\texttt{lore}}: Interface para interagir com a API do lore para buscar dados de patches
    \item \textbf{\texttt{ui}}: Renderização de terminal e lógica de exibição
    \item \textbf{\texttt{cli}}: Parsing de argumentos de linha de comando
    \item \textbf{\texttt{app}}: Estado central da aplicação e gerenciamento de dados
    \begin{itemize}
        \item \texttt{logger}: Gerencia funcionalidade de logging
        \item \texttt{popup}: Cria e gerencia diálogos popup
        \item \texttt{config}: Gerencia configurações da aplicação
    \end{itemize}
    \item \textbf{\texttt{handler}}: Manipulação de entrada do usuário e processamento de eventos
\end{itemize}

\section{Funcionalidades Principais}

A implementação atual fornece várias funcionalidades essenciais para desenvolvedores do kernel:

\begin{itemize}
    \item \textbf{Seleção de Lista de Discussão}: Buscar dinamicamente e navegar listas de discussão arquivadas em lore.kernel.org
    \item \textbf{Últimos Patchsets}: Visualizar os patchsets mais recentes das listas de discussão selecionadas
    \item \textbf{Detalhes e Ações do Patchset}: Visualizar conteúdo individual de patches e metadados (título, autor, versão, número de patches, última atualização, trailers de code-review)
    \item \textbf{Aplicação de Patch}: Aplicar patches a árvores de kernel locais
    \item \textbf{Sistema de Favoritos}: Favoritar patches importantes para referência fácil
    \item \textbf{Integração de Revisão}: Responder com tags \texttt{Reviewed-by} a séries de patches
    \item \textbf{Renderização Aprimorada}: Suporte para ferramentas externas como \texttt{bat}, \texttt{delta} e \texttt{diff-so-fancy} para melhor visualização de patches
\end{itemize}

\section{Limitações Arquiteturais}

Apesar de sua funcionalidade, a implementação atual do Patch-Hub enfrenta desafios arquiteturais significativos que limitam sua flexibilidade e manutenibilidade:

\subsection{Restrições do Sistema de Ownership do Rust}

A limitação mais crítica vem do sistema rigoroso de ownership do Rust. A aplicação usa uma variável central \texttt{app} que mantém configurações, popups e estado da aplicação. Isso cria um problema arquitetural fundamental:

\begin{itemize}
    \item Todo valor em Rust deve ter exatamente um proprietário
    \item Para modificar o app, uma referência mutável (\texttt{\&mut}) é necessária
    \item Popups precisam modificar dados no app, requerendo uma segunda referência \texttt{\&mut} ao app
    \item Rust proíbe ter duas referências \texttt{\&mut} simultâneas aos mesmos dados
\end{itemize}

Esta restrição impede que popups modifiquem dados da aplicação, limitando-os à funcionalidade apenas de exibição. A solução atual envolve refatoração extensiva sempre que novas funcionalidades precisam ser adicionadas.

\subsection{Limitações Single-Threaded}

A aplicação é principalmente single-threaded, o que cria gargalos ao realizar operações de I/O, tais como:

\begin{itemize}
    \item Requisições de rede para a API do lore
    \item Operações do sistema de arquivos
    \item Atualizações de renderização do terminal
    \item Lógica complicada para mostrar telas de carregamento
\end{itemize}

Este design impede que a aplicação lide eficientemente com operações concorrentes e pode levar a interfaces de usuário não responsivas durante operações de rede ou arquivo.

\subsection{Acoplamento Forte}

A arquitetura atual exibe acoplamento forte entre diferentes componentes:

\begin{itemize}
    \item Lógica de UI está entrelaçada com lógica de negócio
    \item Gerenciamento de configuração está profundamente embarcado no estado da aplicação
    \item Tratamento de erro está espalhado por todo o código
    \item Testes se tornam difíceis devido à estrutura monolítica e código com efeitos colaterais usado em todos os lugares
\end{itemize}

\section{Motivação para Redesign}

As limitações descritas acima fornecem um caso convincente para redesign arquitetural. O Modelo de Atores oferece várias vantagens que abordam essas questões:

\subsection{Concorrência e Isolamento}

O Modelo de Atores suporta naturalmente operações concorrentes através de passagem de mensagens, permitindo que a aplicação lide com múltiplas operações de I/O simultaneamente sem bloquear a interface do usuário.

\subsection{Desacoplamento}

Atores se comunicam exclusivamente através de mensagens, criando limites claros entre componentes e reduzindo o acoplamento. Isso torna o sistema mais modular e mais fácil de testar.

\subsection{Isolamento de Estado}

Cada ator gerencia seu próprio estado independentemente, eliminando as questões complexas de ownership que surgem de estado mutável compartilhado em Rust.

\subsection{Extensibilidade}

A arquitetura baseada em mensagens torna mais fácil adicionar novas funcionalidades sem refatoração extensiva, já que novos atores podem ser adicionados sem modificar os existentes.

\section{Objetivos}

Este trabalho não criará uma reescrita completa do Patch-Hub com 100\% de compatibilidade de funcionalidades. Ao invés disso, um subconjunto das funcionalidades foi escolhido para implementação, suficiente para ter as funcionalidades principais do Patch-Hub, que são:

\begin{itemize}
    \item Uma interface de usuário de texto interativa
    \item Integração com API do Lore com suporte a cache
    \item Navegar através das listas de discussão catalogadas
    \item Acesso a um feed de patches enviados para uma lista de discussão
    \item Fornecer uma visualização do conteúdo e metadados de um patch
    \item Configurações personalizáveis
    \item Um sistema de logging
\end{itemize}

\section{Atores Primitivos}

As funcionalidades mencionadas são compostas de algumas capacidades primitivas:

\begin{itemize}
    \item Operações do sistema de arquivos
    \item Requisições HTTP
    \item Manipulação de variáveis de ambiente
    \item Controle de terminal
    \item Integração com shell
\end{itemize}

Capacidades primitivas podem ser entendidas como blocos de construção do sistema que não fazem parte da lógica de negócio e dependem principalmente de interagir com sistemas externos (rede, shell, arquivos, etc). Cada uma dessas capacidades pode ser traduzida para um ator em nossa aplicação seguindo os princípios SOLID, especialmente:

\begin{itemize}
    \item \textbf{Responsabilidade Única}: cada ator lida com um e apenas um domínio
    \item \textbf{Segregação de Interface}: ao dividir em diferentes atores, a injeção de dependência é mais granular, então clientes dependem apenas de interfaces que importam para eles
\end{itemize}

Esses atores serão chamados de atores primitivos, já que não dependem de outros atores.

\subsection{Ator de Sistema de Arquivos}

O ator de Sistema de Arquivos fornece acesso thread-safe a operações do sistema de arquivos através do padrão Modelo de Atores. Ele encapsula todas as operações de arquivo e diretório, fornecendo uma interface unificada tanto para acesso real ao sistema de arquivos quanto para implementações mock para testes.

\textbf{Mensagens:}

\begin{itemize}
    \item \textbf{ReadFile}: Abre um arquivo apenas para leitura (não cria se não existir)
    \item \textbf{WriteFile}: Abre um arquivo para escrita (trunca conteúdo, cria se necessário)
    \item \textbf{AppendFile}: Abre um arquivo para anexar (cria se necessário)
    \item \textbf{RemoveFile}: Remove um arquivo do sistema de arquivos
    \item \textbf{ReadDir}: Lê o conteúdo de um diretório
    \item \textbf{MkDir}: Cria um diretório e seus pais
    \item \textbf{RmDir}: Remove um diretório e seu conteúdo
\end{itemize}

\subsection{Ator de Rede}

O ator de Rede fornece acesso thread-safe a operações HTTP através do padrão Modelo de Atores. Ele gerencia todas as requisições HTTP usando a biblioteca reqwest, fornecendo uma interface unificada tanto para operações reais de rede quanto para implementações mock para testes.

\textbf{Mensagens:}

\begin{itemize}
    \item \textbf{Get}: Executa uma requisição HTTP GET para a URL especificada com headers opcionais
    \item \textbf{Post}: Executa uma requisição HTTP POST para a URL especificada com headers opcionais e corpo
    \item \textbf{Put}: Executa uma requisição HTTP PUT para a URL especificada com headers opcionais e corpo
    \item \textbf{Delete}: Executa uma requisição HTTP DELETE para a URL especificada com headers opcionais
    \item \textbf{Patch}: Executa uma requisição HTTP PATCH para a URL especificada com headers opcionais e corpo
\end{itemize}

\subsection{Ator de Variáveis de Ambiente}

O ator de Ambiente fornece acesso thread-safe a operações de variáveis de ambiente através do padrão Modelo de Atores. Ele encapsula todo o acesso a variáveis de ambiente, fornecendo uma interface unificada tanto para acesso real ao ambiente do sistema quanto para implementações mock para testes.

\textbf{Mensagens:}

\begin{itemize}
    \item \textbf{Set}: Define uma variável de ambiente para um valor especificado
    \item \textbf{Unset}: Remove uma variável de ambiente
    \item \textbf{Get}: Recupera o valor de uma variável de ambiente (retorna VarError se não encontrada)
\end{itemize}

\subsection{Ator de Terminal}

O ator de Terminal fornece acesso thread-safe a operações de UI de terminal através do padrão Modelo de Atores. Ele gerencia o loop de eventos TUI (Terminal User Interface) Cursive e fornece uma interface baseada em mensagens para atualizar a UI de outros atores.

\textbf{Mensagens:}

\begin{itemize}
    \item \textbf{Show}: Atualiza a tela exibida com um novo tipo de tela
    \item \textbf{Quit}: Termina a UI e sai da aplicação
\end{itemize}

\subsection{Ator de Shell}

O ator de Shell fornece acesso thread-safe à execução de programas externos através do padrão Modelo de Atores. Ele encapsula toda a execução de comandos shell, fornecendo uma interface unificada tanto para execução real de processos quanto para implementações mock para testes.

\textbf{Mensagens:}

\begin{itemize}
    \item \textbf{Execute}: Executa um programa externo com comando, argumentos e entrada stdin opcionais dados
\end{itemize}

\section{Outros Atores}

Estes atores se baseiam nos primitivos para implementar a lógica de negócio da aplicação. Eles frequentemente dependem de um ou mais atores primitivos para realizar suas tarefas.

\subsection{Ator de Log}

O ator de Log fornece um serviço de logging centralizado e thread-safe. Ele recebe mensagens de log de outros atores e as escreve em um arquivo de log designado, fornecendo diferentes níveis de log para debugging e monitoramento.

\textbf{Depends on:}

\begin{itemize}
    \item \textbf{Sistema de Arquivos}: Para persistir logs em um arquivo
\end{itemize}

\textbf{Mensagens:}

\begin{itemize}
    \item \textbf{Log}: Registra uma mensagem com um nível de log específico (ex: Info, Warn, Error, Debug)
    \item \textbf{GetLastLogs}: Recupera as últimas N entradas de log
\end{itemize}

\subsection{Ator de Config}

O ator de Config gerencia as configurações da aplicação. Ele é responsável por carregar configurações de um arquivo de configuração, fornecê-las a outros atores sob demanda e persistir mudanças.

\textbf{Depends on:}

\begin{itemize}
    \item \textbf{Sistema de Arquivos}: Para salvar e carregar o arquivo de configuração
\end{itemize}

\textbf{Mensagens:}

\begin{itemize}
    \item \textbf{Get}: Recupera o valor de uma chave de configuração específica
    \item \textbf{Set}: Atualiza o valor de uma chave de configuração específica
    \item \textbf{Load}: Recarrega a configuração do arquivo fonte
    \item \textbf{Persist}: Salva a configuração atual no arquivo fonte
\end{itemize}

\subsection{Ator da API do Lore}

O ator da API do Lore gerencia toda a comunicação com a API externa do Lore patchwork. Ele abstrai os detalhes das requisições HTTP e parsing de dados, fornecendo uma interface limpa e tipada para buscar dados como listas de discussão, séries de patches e patches individuais. Ele atua como um tradutor entre as estruturas de dados internas da aplicação e as respostas brutas da API. Ele não realiza nenhum cache por si só.

\textbf{Depends on:}

\begin{itemize}
    \item \textbf{Rede}: Para executar as requisições HTTP subjacentes
\end{itemize}

\textbf{Mensagens:}

\begin{itemize}
    \item \textbf{GetAvailableListsPage}: Busca uma lista paginada de listas de discussão disponíveis
    \item \textbf{GetPatchFeedPage}: Busca um feed paginado de patches para uma lista de discussão específica
    \item \textbf{GetRawPatch}: Busca o conteúdo bruto de um único patch
\end{itemize}

\subsection{Atores de Cache}

Para atender ao objetivo de fornecer suporte a cache, o sistema usa um conjunto de atores de cache especializados. Ao invés de um cache único e monolítico, a responsabilidade é dividida entre atores, cada um adaptado a um tipo específico de dados e capaz de acessar diretamente a API do Lore para lidar com validação de cache e cache misses. Esta abordagem melhora a modularidade e permite diferentes estratégias de cache para diferentes dados.

\subsubsection{Ator de Cache de Lista de Discussão}

Armazena em cache a lista de todas as listas de discussão disponíveis da API do Lore. Ele busca a lista completa, ordena alfabeticamente e a persiste no sistema de arquivos para acelerar inicializações subsequentes da aplicação. Ele gerencia validação de cache para atualizar periodicamente a lista.

\textbf{Depends on:}

\begin{itemize}
    \item \textbf{API do Lore}: Para buscar os dados da lista de discussão
    \item \textbf{Sistema de Arquivos}: Para persistir o cache em disco
    \item \textbf{Config}: Para obter o caminho do arquivo de cache
    \item \textbf{Log}: Para registrar operações de cache
\end{itemize}

\textbf{Mensagens:}

\begin{itemize}
    \item \textbf{Get}: Recupera uma única lista de discussão por seu índice na lista ordenada
    \item \textbf{GetSlice}: Recupera um intervalo de listas de discussão para paginação
    \item \textbf{Refresh}: Força uma atualização completa do cache da API
    \item \textbf{Invalidate}: Limpa o cache em memória e em disco
    \item \textbf{Len}: Retorna o número total de listas de discussão em cache
\end{itemize}

\subsubsection{Ator de Cache de Feed}

Fornece cache por lista de discussão de metadados de patches (como autor, assunto e data). Ele busca feeds de patches sob demanda, suporta paginação e persiste os metadados no sistema de arquivos. Ele usa uma estratégia de validação inteligente para buscar apenas patches novos quando um feed foi atualizado no servidor.

\textbf{Depends on:}

\begin{itemize}
    \item \textbf{API do Lore}: Para buscar metadados de patches
    \item \textbf{Sistema de Arquivos}: Para persistir o cache em disco
    \item \textbf{Config}: Para obter o caminho do diretório de cache
    \item \textbf{Log}: Para registrar operações de cache
\end{itemize}

\textbf{Mensagens:}

\begin{itemize}
    \item \textbf{Get}: Recupera metadados de um único patch por seu índice dentro de um feed de lista de discussão
    \item \textbf{GetSlice}: Recupera um intervalo de metadados de patches para paginação
    \item \textbf{Refresh}: Atualiza inteligentemente o cache para uma lista de discussão específica, buscando apenas itens novos
    \item \textbf{Invalidate}: Limpa o cache para uma lista de discussão específica
    \item \textbf{Len}: Retorna o número de entradas de metadados em cache para uma lista
\end{itemize}

\subsubsection{Ator de Cache de Patch}

Armazena em cache o conteúdo bruto de patches individuais (formato \texttt{.mbox}). Como o conteúdo do patch é imutável, uma vez que um patch é buscado e armazenado em cache, ele é considerado válido para sempre. Este ator armazena cada patch em um arquivo separado no disco e usa um buffer LRU (Least Recently Used) em memória para fornecer acesso rápido a patches visualizados recentemente.

\textbf{Depends on:}

\begin{itemize}
    \item \textbf{API do Lore}: Para buscar o conteúdo bruto do patch
    \item \textbf{Sistema de Arquivos}: Para armazenar arquivos de patch permanentemente
    \item \textbf{Config}: Para obter o caminho do diretório de cache
    \item \textbf{Log}: Para registrar operações de cache
\end{itemize}

\textbf{Mensagens:}

\begin{itemize}
    \item \textbf{Get}: Recupera o conteúdo bruto de um patch dado sua lista de discussão e ID de mensagem
    \item \textbf{Invalidate}: Remove um patch específico do cache em disco e buffer em memória
    \item \textbf{IsAvailable}: Verifica se um patch está em cache sem buscá-lo
\end{itemize}

\section{Métricas}

% PLACEHOLDER: This section is marked as incomplete in Index.md
% TODO: Add content about metrics (feature completeness, performance, test coverage)
% This section should cover:
% - Feature completeness comparison between original and actor-based implementation
% - Performance benchmarks and analysis
% - Test coverage metrics
% - Qualitative analysis of maintainability, testability, extensibility
% - Comparison of architectural benefits and trade-offs

\subsection{Completude de Funcionalidades}

% PLACEHOLDER: Content to be developed
% Compare the feature set of the original Patch-Hub with the actor-based implementation
% Identify which features were successfully implemented and which were omitted
% Discuss the rationale for feature selection

\subsection{Performance}

% PLACEHOLDER: Content to be developed
% Present benchmarks comparing performance metrics between implementations
% Analyze memory usage, CPU utilization, and response times
% Discuss the impact of the actor model on performance characteristics

\subsection{Cobertura de Testes}

% PLACEHOLDER: Content to be developed
% Present test coverage metrics for the actor-based implementation
% Compare testing approaches between the two architectures
% Discuss the benefits of actor-based testing strategies

\subsection{Análise Qualitativa}

% PLACEHOLDER: Content to be developed
% Evaluate maintainability improvements
% Assess testability enhancements
% Analyze extensibility benefits
% Discuss developer experience and code organization

